\begin{figure}[H]
  \centering
  \begin{minipage}[c]{0.42\textwidth}
    \centering
    \begin{tikzpicture}[scale=0.92]
      \node[graph vertex] (v1) at (0,2.2) {$v_1$};
      \node[graph vertex] (v2) at (1.8,1.2) {$v_2$};
      \node[graph vertex] (v3) at (-1.8,1.2) {$v_3$};
      \node[graph vertex] (v4) at (0,0.2) {$v_4$};
      \node[graph vertex] (v5) at (1.8,-0.8) {$v_5$};
      \node[graph vertex] (v6) at (1.8,-2.0) {$v_6$};
      \draw[graph edge] (v1)--(v2) (v1)--(v3) (v2)--(v3);
      \draw[graph edge] (v2)--(v4) (v3)--(v4) (v4)--(v5) (v5)--(v6);
      \draw[graph edge] (v4) to[out=270,in=180,looseness=8] (v4);
    \end{tikzpicture}
    \\[2mm]
    \small Graf nieskierowany
  \end{minipage}
  \hfill
  \begin{minipage}[c]{0.52\textwidth}
    \small
    \begin{tabular}{@{}r@{\;}l@{}}
      $v_1 \to$ & $[v_2, v_3]$ \\
      $v_2 \to$ & $[v_1, v_3, v_4]$ \\
      $v_3 \to$ & $[v_1, v_2, v_4]$ \\
      $v_4 \to$ & $[v_2, v_3, v_4, v_5]$ \\
      $v_5 \to$ & $[v_4, v_6]$ \\
      $v_6 \to$ & $[v_5]$
    \end{tabular}
    \\[2mm]
    \small Lista sąsiedztwa
  \end{minipage}
  \caption{Lista sąsiedztwa przykładowego grafu nieskierowanego o sześciu wierzchołkach}
  \label{fig:lista-sasiedztwa-graf-nieskierowany}
\end{figure}
